\chapter{Bloc 3 --- Entrepôt OMEGA BI (C13-C17)}
\label{chap:bloc3}

Le Bloc 3 construit et maintient l'entrepôt décisionnel OMEGA BI, qui
répond au besoin d'analyse de la performance logistique exprimé par la
direction. Clos le 21/09/2026 (milestone M2), les 5 compétences ont
chacune un livrable vérifié -- détail exhaustif dans
\filepath{docs/comptes_rendus/M2.md}.

\section{Séquencement retenu}

L'ordre C13~$\rightarrow$~C14~$\rightarrow$~C15~$\rightarrow$~C17~$\rightarrow$~C16
(\filepath{docs/architecture/sequencement_bloc3.md}) inverse volontairement
C16/C17 par rapport au Gantt initial : la procédure de purge
\hyperref[gloss:rgpd]{RGPD} de C16 dépend du schéma historisé que C17
doit d'abord créer.

\section{C13 --- Modéliser un entrepôt de données}

Approche bottom-up par \hyperref[gloss:datamart]{datamarts} (Kimball) :
2 datamarts, 3 faits, 6 dimensions conformées, matrice de bus
(\filepath{docs/architecture/modelisation_omega_bi.md}). Le seul schéma
en flocon retenu (\texttt{Dim\_Produit}~$\rightarrow$~\texttt{Dim\_Categorie})
est justifié par une contrainte réelle des sources, pas par principe de
modélisation.

\begin{figure}[htbp]
  \centering
  \includegraphics[width=\textwidth]{modelisation_omega_bi_mcd}
  \caption{Modèle conceptuel de données de l'entrepôt OMEGA BI (\filepath{docs/architecture/modelisation_omega_bi.md})}
\end{figure}

\textbf{Pourquoi pas de \texttt{Fait\_Livraison} séparée} :
\texttt{livraisons} (staging, TransFlow) et \texttt{expeditions}
(FluxPro) décrivent le même événement métier vu par deux systèmes
sources, rapprochés par \texttt{tracking\_number}
(chapitre~\ref{chap:bloc2} §C11). La cardinalité réelle a été vérifiée
avant de modéliser, pas supposée : 100~\% des 1100 expéditions FluxPro
ont exactement une livraison TransFlow correspondante (relation 1:1
confirmée empiriquement). Scinder deux sources en cardinalité 1:1
vérifiée en deux tables de faits aurait imposé une jointure
fait-à-fait à chaque requête -- anti-pattern Kimball qu'une dimension
conformée évite normalement. Les mesures de \texttt{livraisons}
(délai, statut) sont donc absorbées directement dans
\texttt{Fait\_Expedition}, au même grain.

\section{C14 --- Créer un entrepôt de données}

L'entrepôt est créé dans une base PostgreSQL distincte de la base de
staging, mais sur la même instance (sobriété \hyperref[gloss:rgesn]{RGESN}) -- 3 schémas
physiques créés à ce stade (\texttt{dimensions}, \texttt{exploitation},
\texttt{commercial}), versionnés via un second environnement Alembic
dédié. Le cycle \texttt{upgrade}/\texttt{downgrade} est vérifié contre une
base réelle, pas seulement contre une base vide (un bug latent
d'Alembic, \texttt{include\_schemas=True} manquant, n'a d'ailleurs été
révélé qu'à la seconde migration). Deux schémas supplémentaires
rejoignent la même base par la suite, sans modifier les trois premiers
: \texttt{gouvernance} (journalisation des opérations, C16) et
\texttt{monitoring} (métriques du data lake, jalon M4/C20bis,
chapitre~\ref{chap:m4}).

\begin{figure}[htbp]
  \centering
  \includegraphics[width=\textwidth]{creation_entrepot_omega_bi_schemas}
  \caption{Les 5 schémas PostgreSQL de \texttt{datacore\_omega\_bi} et leurs liens (\filepath{docs/architecture/creation_entrepot_omega_bi.md})}
\end{figure}

Seules les tables de faits référencent les dimensions par clé
étrangère inter-schéma ; \texttt{gouvernance} et \texttt{monitoring}
sont des journaux techniques autonomes, jamais joints aux dimensions.
\texttt{bi\_reader} (lecture seule) a accès à tous les schémas sauf
\texttt{gouvernance}, délibérément exclu (préoccupation d'exploitation
interne, pas de restitution BI).

\section{C15 --- Développer des pipelines de données}

Le pipeline \hyperref[gloss:etl]{ETL} staging~$\rightarrow$~entrepôt applique des contrôles
qualité réels (mise en quarantaine des lignes non rapprochées de
l'historique, pas un rejet silencieux), et sa cohérence est vérifiée
en comparant les indicateurs calculés par l'ETL à ceux calculés
indépendamment par l'API C12 -- deux chemins de calcul distincts
produisant le même résultat.

\section{C16 --- Assurer la gestion opérationnelle d'un environnement de données}

La gestion opérationnelle (\filepath{docs/architecture/gestion_operationnelle_omega_bi.md})
couvre des sauvegardes planifiées (partielle/complète, justifiées par
la reconstructibilité de 7 des 8 tables depuis le staging), une
journalisation réelle des opérations dans le nouveau schéma
\texttt{gouvernance} (succès et échec, sur une connexion dédiée pour
survivre à un \texttt{ROLLBACK}), et un registre RGPD appuyé sur un
audit automatisé répétable ne trouvant aucune donnée personnelle dans
l'entrepôt.

Ce périmètre a été étendu après la clôture initiale du bloc : la
restitution par vues SQL et notebook s'est révélée insuffisante comme
preuve de restitution, et renforcée dans un outil de
\hyperref[gloss:bi]{BI} dédié -- traité comme son propre jalon
transverse (chapitre~\ref{chap:m4}).

\section{C17 --- Mettre en œuvre une variation de dimension}

\texttt{Dim\_Client} est historisée en \hyperref[gloss:scd2]{SCD2}
(type~2 de Kimball, voir aussi \hyperref[gloss:scd1]{SCD1}) : un
changement de \texttt{nom} ou \texttt{secteur} ferme la version
courante et en ouvre une nouvelle, plutôt que d'écraser la valeur
précédente.

\begin{lstlisting}[language=Python,caption={Détection du changement et historisation SCD2 (\filepath{src/datacore/storage/warehouse/load_warehouse.py})}]
if current is not None and current[1] == r["nom"] and current[2] == r["secteur"]:
    mapping[r["id"]] = current[0]
    continue

if current is not None:
    cur.execute(
        """
        UPDATE dimensions.dim_client
        SET valid_to = %(today)s, is_current = false
        WHERE client_key = %(client_key)s
        """,
        {"today": today, "client_key": current[0]},
    )
# ... puis INSERT d'une nouvelle ligne avec une nouvelle client_key.
\end{lstlisting}

\texttt{r} est la ligne source lue depuis le staging (dict avec
\texttt{nom}/\texttt{secteur}/\texttt{id} FluxPro) ; \texttt{current}
est la version actuellement courante de ce client dans l'entrepôt si
elle existe déjà (\texttt{None} sinon), sous forme de tuple
\texttt{(client\_key, nom, secteur)} ; \texttt{mapping} associe
l'identifiant FluxPro à la \texttt{client\_key} à utiliser pour charger
les faits qui le référencent, qu'une nouvelle version ait été créée ou
non. Le mécanisme est vérifié par un changement réel simulé : l'ancienne
version reste close et interrogeable, la nouvelle est créée, et les
faits déjà chargés continuent de référencer la version qui était
courante au moment de leur chargement. Écart assumé avec la
description initiale de l'issue (qui demandait l'historisation
d'attributs absents du jeu de données pédagogique, adresse et
contrat) : \texttt{nom} et \texttt{secteur} ont été retenus à la
place, après décision prise avant implémentation plutôt que devinée.
